\documentclass[11pt,a4paper]{ctexart}
\usepackage[margin=1.6cm]{geometry}
\usepackage{amsmath,amssymb}
\usepackage{booktabs,array}
\usepackage{xcolor}
\usepackage{tcolorbox}
\setlength{\parindent}{0pt}
\setlength{\parskip}{5pt}
\renewcommand{\arraystretch}{1.35}

\title{\vspace{-1.4em}粘性流动作业-2：后续题完整答案汇总}
\author{}
\date{}

\begin{document}
\maketitle
\vspace{-2.2em}

\begin{tcolorbox}[colback=blue!3,colframe=blue!55!black,title=范围]
本文件汇总后续题：书 10.3、书 10.5、书 10.10、第8题。
\end{tcolorbox}

\section*{一、书 10.3：排挤厚度和动量损失厚度}

定义：
\[
\delta_1=\int_0^\infty\left(1-\frac{u}{U}\right)\,dy,\qquad
\delta_2=\int_0^\infty\frac{u}{U}\left(1-\frac{u}{U}\right)\,dy.
\]

由于 \(y>\delta\) 时 \(u=U\)，只需从 \(0\) 积到 \(\delta\)。令
\[
\eta=\frac{y}{\delta},\qquad dy=\delta d\eta.
\]

\subsection*{(1) \(u/U=\eta\)}

\[
\delta_1=\delta\int_0^1(1-\eta)d\eta=\frac{\delta}{2}.
\]

\[
\delta_2=\delta\int_0^1\eta(1-\eta)d\eta=\frac{\delta}{6}.
\]

\[
\boxed{\delta_1=\frac{\delta}{2},\qquad \delta_2=\frac{\delta}{6}}
\]

\subsection*{(2) \(u/U=\frac32\eta-\frac12\eta^3\)}

记
\[
f(\eta)=\frac32\eta-\frac12\eta^3.
\]

\[
\delta_1=\delta\int_0^1(1-f)d\eta
=\delta\int_0^1\left(1-\frac32\eta+\frac12\eta^3\right)d\eta
=\frac38\delta.
\]

\[
\delta_2=\delta\int_0^1f(1-f)d\eta
=\delta\left(\int_0^1f\,d\eta-\int_0^1f^2\,d\eta\right).
\]

\[
\int_0^1f\,d\eta=\frac58,\qquad
\int_0^1f^2\,d\eta=\frac{17}{35}.
\]

\[
\delta_2=\delta\left(\frac58-\frac{17}{35}\right)=\frac{39}{280}\delta.
\]

\[
\boxed{\delta_1=\frac38\delta,\qquad \delta_2=\frac{39}{280}\delta}
\]

\section*{二、书 10.5：正弦速度分布动量积分}

设
\[
\frac{u}{U_e}=a+b\sin\left(\frac{cy}{\delta}\right).
\]

边界条件：
\[
y=0:\quad u=0\Rightarrow a=0.
\]

\[
y=\delta:\quad u=U_e\Rightarrow b\sin c=1.
\]

\[
y=\delta:\quad \frac{\partial u}{\partial y}=0
\Rightarrow b\frac{c}{\delta}\cos c=0
\Rightarrow c=\frac{\pi}{2},\quad b=1.
\]

所以
\[
\frac{u}{U_e}=\sin\left(\frac{\pi y}{2\delta}\right).
\]

令 \(\eta=y/\delta\)，
\[
f(\eta)=\sin\frac{\pi\eta}{2}.
\]

排挤厚度：
\[
\delta_1=\delta\int_0^1\left(1-\sin\frac{\pi\eta}{2}\right)d\eta
=\delta\left(1-\frac{2}{\pi}\right).
\]

动量损失厚度：
\[
\delta_2=\delta\int_0^1 f(1-f)d\eta
=\delta\left(\frac{2}{\pi}-\frac12\right).
\]

因此：
\[
\frac{\delta_1}{\delta_2}
=\frac{1-\frac{2}{\pi}}{\frac{2}{\pi}-\frac12}
\approx2.66.
\]

\[
\boxed{\frac{\delta_1}{\delta_2}=2.66}
\]

零压强梯度平板动量积分关系式：
\[
\frac{\tau_w}{\rho U_e^2}=\frac{d\delta_2}{dx}.
\]

壁面切应力：
\[
\tau_w=\mu\left(\frac{\partial u}{\partial y}\right)_{0}
=\mu U_e\frac{\pi}{2\delta}.
\]

设
\[
\delta_2=A\delta,\qquad A=\frac{2}{\pi}-\frac12.
\]

则
\[
A\frac{d\delta}{dx}=\frac{\nu\pi}{2U_e\delta}.
\]

积分得
\[
\delta^2=\frac{\pi\nu x}{AU_e}.
\]

平板单面总阻力系数：
\[
C_D=\frac{2\delta_2(L)}{L}
=2\sqrt{\frac{A\pi}{Re_L}}.
\]

因为
\[
2\sqrt{A\pi}=1.31,
\]

所以
\[
\boxed{C_D=\frac{1.31}{\sqrt{Re_L}}}
\]

\section*{三、书 10.10：流线型火车摩擦阻力功率}

已知：
\[
L=110\ \mathrm{m},\quad b=8.25\ \mathrm{m},\quad
S=Lb=907.5\ \mathrm{m^2}.
\]

\[
U=160\ \mathrm{km/h}=44.44\ \mathrm{m/s}.
\]

\[
\rho=1.22\ \mathrm{kg/m^3},\qquad
\mu=1.79\times10^{-5}\ \mathrm{Pa\cdot s}.
\]

临界雷诺数：
\[
Re_{cr}=5\times10^5.
\]

\subsection*{1. 层流边界层维持长度}

\[
Re_{cr}=\frac{\rho Ux_{cr}}{\mu}.
\]

\[
x_{cr}=\frac{Re_{cr}\mu}{\rho U}
=\frac{5\times10^5\times1.79\times10^{-5}}
{1.22\times44.44}.
\]

\[
\boxed{x_{cr}\approx0.165\ \mathrm{m}}
\]

\subsection*{2. 摩擦功率}

\[
Re_L=\frac{\rho UL}{\mu}
=\frac{1.22\times44.44\times110}{1.79\times10^{-5}}
\approx3.33\times10^8.
\]

有层流到湍流转捩时，光滑平板平均摩擦阻力系数取
\[
\bar C_f=\frac{0.074}{Re_L^{1/5}}-\frac{1742}{Re_L}.
\]

\[
\bar C_f\approx0.001456.
\]

阻力：
\[
D=\bar C_f\frac12\rho U^2S.
\]

\[
D=0.001456\times\frac12\times1.22\times44.44^2\times907.5
\approx1.59\times10^3\ \mathrm{N}.
\]

功率：
\[
P=DU.
\]

\[
P=1.59\times10^3\times44.44\approx7.08\times10^4\ \mathrm{W}.
\]

\[
\boxed{P\approx70.8\ \mathrm{kW}}
\]

\section*{四、第8题：不同流体绕平板比较}

已知：
\[
U=2.5\ \mathrm{m/s},\quad L=0.5\ \mathrm{m},\quad b=3\ \mathrm{m}.
\]

\[
S=Lb=1.5\ \mathrm{m^2},\qquad Re_{cr}=5\times10^5.
\]

\subsection*{(a) 空气}

\[
\rho=1.2\ \mathrm{kg/m^3},\qquad
\nu=1.5\times10^{-5}\ \mathrm{m^2/s}.
\]

\[
Re_L=\frac{UL}{\nu}
=\frac{2.5\times0.5}{1.5\times10^{-5}}
=8.33\times10^4<Re_{cr}.
\]

所以为空气全层流。

\[
\bar C_f=\frac{1.328}{\sqrt{Re_L}}
=0.00460.
\]

\[
D=\bar C_f\frac12\rho U^2S
=0.00460\times\frac12\times1.2\times2.5^2\times1.5.
\]

\[
\boxed{D_{\mathrm{air}}\approx0.0259\ \mathrm{N}}
\]

尾部边界层厚度：
\[
\delta_L=\frac{5L}{\sqrt{Re_L}}
=\frac{5\times0.5}{\sqrt{8.33\times10^4}}
\approx8.66\times10^{-3}\ \mathrm{m}.
\]

\[
\boxed{\delta_{\mathrm{air}}\approx8.66\ \mathrm{mm}}
\]

\subsection*{(b) 水}

\[
\rho=998\ \mathrm{kg/m^3},\qquad
\nu=1.005\times10^{-6}\ \mathrm{m^2/s}.
\]

\[
Re_L=\frac{UL}{\nu}
=\frac{2.5\times0.5}{1.005\times10^{-6}}
=1.24\times10^6>Re_{cr}.
\]

会发生转捩。

\[
x_{cr}=\frac{Re_{cr}\nu}{U}
=\frac{5\times10^5\times1.005\times10^{-6}}{2.5}
=0.201\ \mathrm{m}.
\]

平均摩擦阻力系数：
\[
\bar C_f=\frac{0.074}{Re_L^{1/5}}-\frac{1742}{Re_L}
\approx0.00307.
\]

\[
D=\bar C_f\frac12\rho U^2S.
\]

\[
D=0.00307\times\frac12\times998\times2.5^2\times1.5.
\]

\[
\boxed{D_{\mathrm{water}}\approx14.4\ \mathrm{N}}
\]

尾部为湍流边界层，名义厚度取：
\[
\delta_L=\frac{0.37L}{Re_L^{1/5}}.
\]

\[
\delta_L=\frac{0.37\times0.5}{(1.24\times10^6)^{1/5}}
\approx1.12\times10^{-2}\ \mathrm{m}.
\]

\[
\boxed{\delta_{\mathrm{water}}\approx11.2\ \mathrm{mm}}
\]

\subsection*{比较}

\begin{center}
\begin{tabular}{lllll}
\toprule
流体 & \(Re_L\) & 状态 & 单面摩擦阻力 & 尾部边界层厚度\\
\midrule
空气 & \(8.33\times10^4\) & 全层流 & \(0.0259\ \mathrm{N}\) & \(8.66\ \mathrm{mm}\)\\
水 & \(1.24\times10^6\) & 转捩后湍流 & \(14.4\ \mathrm{N}\) & \(11.2\ \mathrm{mm}\)\\
\bottomrule
\end{tabular}
\end{center}

水的运动粘度小，雷诺数大，容易转捩为湍流；同时水的密度远大于空气，所以摩擦阻力大很多。

\end{document}
